\documentclass[11pt]{article}
\usepackage[margin=1in]{geometry}
\usepackage{cmap}
\usepackage[T1]{fontenc}
\input{glyphtounicode}
\pdfgentounicode=1
\usepackage{booktabs,hyperref,amsmath}
\hypersetup{colorlinks=true,linkcolor=blue,urlcolor=blue}
\sloppy
\title{\textbf{Response: five items, and a connection drop}}
\author{}\date{11 October 2026}
\begin{document}\maketitle\vspace{-3em}

All five applied. One of them came back against the premise of the objection, and your
machine disconnected from this session partway through, so the last section lists what is
not yet on your disk.

\section*{1. API selection rule --- measured, and there are no duplicates}

This was answerable by measurement rather than by a statement, so I measured it. The cached
responses contain zero duplicate \texttt{(material\_id, thermo\_type)} combinations:
152{,}426 unique identifiers of 152{,}426 rows for \texttt{GGA\_GGA+U}, 37{,}123 of
37{,}123 for r$^2$SCAN, and 22{,}756 of 22{,}756 in the completeness pull. The pull script
performs no deduplication of its own, so the API returned exactly one document per
material per thermo type.

Appendix E now says that, and says what it means: the dictionary construction would retain
the last document if a material ever yielded two of the same type, but no selection rule is
exercised in this work, and a replication that does hit duplicates should not rely on
insertion order.

\section*{2 and 3. Abstract wording --- your phrasings, verbatim}

``Presence of a Materials Project Hubbard-$U$-list element is strongly conditionally
associated with disagreement in the fitted model'' and ``Neither selective $U$ application
nor hull proximity explains away this association under the analyses performed''. The
``All'' sentence now reads ``Under the release documentation's interpretation of its `All'
reference set, the two cells sit at 14.5\% and 37.4\%, preserving a 22.9-point spread while
substantially reducing the ratio.''

The abstract is 287 words. Your 250 target and the three clauses you have required of it
are not simultaneously satisfiable; the hedging costs words. I am not going to quietly pick
which of your requirements to drop.

\section*{5. ``Leading candidate'' --- gone, in both places}

Both the Conclusion and \S4.3 now say ``a plausible candidate'', and both add that no
systematic comparison of candidate mechanisms was run and none are ranked. The phrase was
mine, taken from the fix list's ``leading candidate'' wording, and your objection to it is
right: nothing in the paper establishes a ranking.

\section*{4. Table 4 --- the premise was wrong, and the uncertainty is now there}

Two separate things, and the first is a correction to the objection.

\paragraph{The design is not rank-deficient.} The objection says the model ``uses a
minimum-norm solution for residual rank deficiency''. It does not, because there is no
rank deficiency. Measured: 113 columns, 83 element fractions and 30 surviving interactions,
numerical rank 113 at a tolerance of $1.6\times10^{-9}$, empty null space, condition number
$7.0\times10^{2}$. The element-fraction columns do sum to one in every row, but the design
carries no intercept column, so that creates no dependency. The least-squares solution is
unique, every within-element contrast is identified, and the projection of each of the 75
such contrasts onto the numerical null space is exactly zero because the null space is
empty.

That also corrects the paper's own earlier claim. Draft 27's caption said ``the row spread
is identified while the level is not''. That was wrong as a statement about rank: both are
estimated. The caption and text now give the real reason the levels should not be read as
physical contributions, which is that $\Delta$ contains the correction removal by
construction.

\paragraph{The spreads now carry cluster-aware intervals.} Cluster bootstrap over
structural prototypes, 400 replicates, refitting the model on each resampled draw:

\begin{center}
\begin{tabular}{lrl}
\toprule
Element & Spread (meV/atom) & 95\% CI \\
\midrule
Mn & 273 & [196, 342] \\
Cr & 249 & [144, 399] \\
Fe & 239 & [157, 328] \\
Ni & \phantom{0}92 & [47, 212] \\
Co & \phantom{0}36 & [22, 99] \\
\bottomrule
\end{tabular}
\end{center}

All five exclude zero, so the context dependence is not a fitting artefact. The
\emph{ordering} among them is a different question and the text now says so: the Co and Ni
intervals overlap each other, and Ni's overlaps the lower end of Fe's. The replicate count
is 400 rather than the tens of thousands used elsewhere because each replicate refits a
113-column model; it is stated in the caption with the reason.

\paragraph{Two bugs in my own checking, caught by the check.} The first run asserted that
each contrast is invariant under an equivalent reparameterisation and failed, reporting a
2.0 meV/atom discrepancy. The model was fine; my test was not. Dropping one interaction
column per element is not an invertible reparameterisation --- with a full-rank design it
fits a restricted model, not the same model in new coordinates. I replaced it with column
scaling and column permutation, which are invertible. That failed too, at 250 meV/atom,
because I mapped the scaled fit back wrongly: for $XD\beta' = y$ the original-basis
coefficients are $D\beta'$, not $\beta'/d$. The corrected line is written but was not run,
for the reason in the next section.

The identification claim in the paper does not depend on that test. Full rank implies a
unique least-squares solution, which implies every linear functional of it is invariant
under any invertible reparameterisation. The empirical test is belt and braces on a
theorem, and rank was measured.

\section*{What is not on your machine}

Your computer dropped off this session's bridge during the last call, so the following
exist only here and need to be put back:

\begin{itemize}
\item \texttt{paper\_draft28.tex} and \texttt{.pdf} --- everything above. Your
\texttt{paper/gnome\_audit\_paper.tex} is still Draft 27.
\item The corrected scaling line in \texttt{phase43\_context\_identify.py}. The copy on
your disk still has \texttt{c\_scaled/d} and so still exits 1 on its own invariance check.
Its two output CSVs were written by the run that produced the numbers above, so
\texttt{results/context\_spread\_ci.csv} and \texttt{results/context\_identification.csv}
are present and correct apart from the reparameterisation row.
\item \texttt{phase43} is not in \texttt{run\_all.py}'s \texttt{ORDER} and has no headline
checks; the full-rank assertion and the five intervals are therefore not yet covered by
layer one.
\item \texttt{check\_paper.py} does not yet parse Table 4's new CI column.
\item The duplicate check is stated in the paper but not asserted in
\texttt{q2\_known\_baseline.py}. It should be, so a future pull that does return duplicates
fails loudly instead of silently keeping the last row.
\item None of the three checkers has run against Draft 28.
\end{itemize}

\paragraph{One discrepancy to reconcile when they do.} \texttt{context\_spread\_ci.csv}
rounds nickel's spread to 93, computed as max minus min of the unrounded effective
coefficients. Table 4 shows 92, which is max minus min of the \emph{rounded} per-context
values, $+84 - (-8)$. Both are defensible and they disagree by one. The fix is for
\texttt{phase43} to read the point spreads from \texttt{context\_dependence.csv} rather
than recompute them, so the two files cannot diverge. Until that is done, do not let
\texttt{check\_paper.py} compare the Table 4 spread column against
\texttt{context\_spread\_ci.csv}; it will fail on nickel, correctly.

Reconnect the machine and I will commit Draft 28, apply the two code fixes, wire
\texttt{phase43} into the driver with the full-rank assertion as an \texttt{[impl]} check,
extend the table parser, and run all three checkers.

\end{document}
